% ============================
% ESTILO DE PÁGINA
% ============================
%
% CONTENIDO:
%   1. Encabezado/pie   — lhead (sección), rhead (página), cfoot, headrulewidth
%   2. Listados         — firstnumber, numberfirstline (extiende lstset de Preambulo)
%   3. Opcional         — marco de página con fancypage (desactivado)
%
% ============================

% --- Encabezado y pie de página ---
% NOTA: \pagestyle{fancy} debe activarse en el documento principal (main.tex)
%       Se mantiene comentado aquí para evitar conflictos con la carátula.
%\pagestyle{fancy}
\lhead{\tiny\rightmark}            % Sección actual (izquierda)
\rhead{\thepage}                   % Número de página (derecha)
\cfoot{}                           % Pie de página central vacío
\renewcommand{\headrulewidth}{1pt} % Línea decorativa bajo el encabezado

% --- Extensión del lstset global (definido en Preambulo.tex) ---
% Solo se agregan campos no presentes en el lstset base.
\lstset{
  firstnumber=1,        % Reiniciar numeración desde 1 en cada listado
  numberfirstline=true  % Numerar también la primera línea
}

% ============================
% OPCIONAL: MARCO EN LAS PÁGINAS
% ============================
% Descomenta si querés un borde visible en cada página.
% Requiere \usepackage{fancyhdr,geometry,fancybox}

% \fancypage{\setlength{\fboxsep}{10pt}\fbox}{}
% - Dibuja un marco rectangular alrededor del área de texto
% - Ajustá \fboxsep para cambiar la distancia entre texto y borde
